% Part: many-valued-logic
% Chapter: syntax-and-semantics
% Section: formulas

\documentclass[../../../include/open-logic-section]{subfiles}

\begin{document}

\olfileid[fr]{mvl}{syn}{fml}

\olsection{Formules}

\begin{defn}[Formule]
\ollabel{defn:formulas}
L'ensemble~$\Frm[L]$ des
\emph{formules} d'un langage
propositionnel~$\Lang L$ est défini
par induction comme suit :
\begin{enumerate}
\item Toute variable propositionnelle
  $\Obj p_i$ est une formule atomique.
\item Tout connecteur d'arité~$0$
  (constante propositionnelle) de
  $\Lang L$ est une formule atomique.
\item Si $\star$ est un connecteur d'arité~$n\geq1$ de $\Lang L$
  et si $!A_1,\dots,!A_n$ sont des formules, alors
  $\star(!A_1,\dots,!A_n)$ est une formule.
\tagitem{limitClause}{Rien d'autre
  n'est une formule.}{}
\end{enumerate}
Si $\star$ est unaire, nous écrirons
souvent $\star !A_1$ au lieu de
$\star(!A_1)$. S'il est binaire,
nous écrirons souvent
$(!A_1\star !A_2)$ au lieu de
$\star(!A_1,!A_2)$.
\end{defn}

Comme d'habitude, nous omettrons
souvent les parenthèses extérieures.

\begin{ex}
Dans le langage usuel~$\Lang{L_0}$,
$\Obj p_1\lif(\Obj p_1\land
\lnot\Obj p_2)$ est une formule.
Dans le langage de la logique produit,
on l'écrirait plutôt
$\Obj p_1\lif(\Obj p_1\odot
\lnot\Obj p_2)$. Si nous ajoutons
au langage le connecteur unaire
$\triangle$, nous obtenons aussi
des formules telles que
$\triangle(\Obj p_1\land\Obj p_2)
\lif(\triangle\Obj p_1\land
\triangle\Obj p_2)$.
\end{ex}

\end{document}
